\subsection{Support operations in the logarithmic form topology}
\label{sec:window-support-topology}

The cofinal implication holds on the original form domain.  We now
make its approximation maps explicit.  This also permits elementary
vectors with nonzero endpoint values to detect changes in the form,
and then transfers those differences to compact smooth tests.

Write $\widehat g$ for the fixed Fourier transform, always taken
after zero extension outside the specified support.  Put
\[
 w(\xi)=m(\xi)-m_0+1\ge1,\qquad
 \|g\|_X^2=\frac1{2\pi}\int_{\mathbb R}
                 w(\xi)|\widehat g(\xi)|^2\,d\xi.
\]
The symbol estimates give
$w(\xi)\asymp\log(2+|\xi|)$.  For $I_T=(-T,T)$, denote by
$X_T$ the closed support subspace with this norm.  It is exactly
the original logarithmic form domain $\mathcal Q_T$ equipped
with an equivalent Hilbert form norm.  No value or trace is
prescribed at either endpoint.
More explicitly, if $B_T$ is the bounded prime-plus-polar
operator and $\beta_T\ge\|B_T\|$, then
\[
 \|g\|_X^2\le
 Q_T(g)+(1-m_0+\beta_T)\|g\|_2^2
 \le(1+2\beta_T)\|g\|_X^2.
\]
Thus the norm used below is precisely equivalent to a shifted
closed-form norm of the complete operator.

Reflection and conjugation preserve the even Fourier weight, so
$g\mapsto g^*$ is an antilinear isometry of $X_T$.
Write $\mathfrak q_T$ for the polarization of $Q_T$.  It is
continuous on $X_T$.  In fact, if
\[
 C_T=1+|m_0-1|+2T+2\sinh T
              +2\sum_{2\le n\le e^{2T}}\frac{\Lambda_{\mathrm{vM}}(n)}{\sqrt n},
\]
then
\[
 |\mathfrak q_T(f,g)|\le C_T\|f\|_X\|g\|_X,
\]
and consequently
\begin{equation}
 |Q_T(f)-Q_T(g)|
 \le C_T(\|f\|_X+\|g\|_X)\|f-g\|_X.
 \label{eq:window-support-1}
\end{equation}
Indeed the archimedean form is the $X$ inner product plus
$(m_0-1)$ times the $L^2$ inner product.  The remaining terms
are the actual finite prime-power operator and the actual
rank-two polar operator, whose norms have the stated bounds.

\subsubsection{Zero extension and inward dilation}

For $U\ge T$, zero extension gives an isometry
\[
 X_T\longrightarrow X_U
\]
for the displayed norm.  The Fourier transform does not change.
The complete forms agree on this inclusion:
\begin{equation}
 Q_U(g)=Q_T(g)\qquad(g\in X_T).
 \label{eq:window-support-2}
\end{equation}
The polar integrals and archimedean integral are unchanged, and
every new shift $b\ge2T$ has zero correlation with a function
supported in $[-T,T]$.  The equality case concerns only a set
of measure zero and also gives zero.

For $0<\theta\le1$, define the unitary $L^2$ dilation
\[
 (U_\theta g)(x)=\theta^{-1/2}g(x/\theta),\qquad
 \widehat{U_\theta g}(\xi)=\theta^{1/2}\widehat g(\theta\xi).
\]
It maps $X_T$ into $X_{\theta T}$.  For every fixed
$\theta_0>0$ there is a finite constant $A_{\theta_0}$ such
that
\[
 w(\xi/\theta)\le A_{\theta_0}w(\xi)
       \qquad(\theta_0\le\theta\le1).
\]
This follows from equivalence with the logarithmic weight and
\[
 \log(2+|\xi|/\theta)
 \le\log(2+|\xi|)+\log(1/\theta_0).
\]
Thus $U_\theta$ is uniformly bounded on the full weighted
Fourier space when $\theta\in[\theta_0,1]$.
Its strong continuity at one follows first for smooth Fourier
functions of compact support, by change of variables and
dominated convergence, and then for all weighted $L^2$
functions by density and the uniform operator bound.  In
particular
\begin{equation}
 \|U_\theta g-g\|_X\longrightarrow0\qquad(\theta\uparrow1).
 \label{eq:window-support-3}
\end{equation}
This proof does not use a classical endpoint condition.

\subsubsection{Convolution and smooth approximation}

For $f\in X_T$ and $g\in X_U$, their compact supports imply
$f,g\in L^1$, and
\[
 \operatorname{supp}(f\star g)\subset[-T-U,T+U],\qquad
 \widehat{f\star g}=\widehat f\widehat g.
\]
Therefore
\begin{equation}
 \|f\star g\|_X
 \le\min\{\|f\|_1\|g\|_X,\ \|g\|_1\|f\|_X\}.
 \label{eq:window-support-4}
\end{equation}
Since $\|f\|_1\le\sqrt{2T}\|f\|_2\le\sqrt{2T}\|f\|_X$,
this proves joint continuity
$X_T\times X_U\to X_{T+U}$.  In particular convergence of
both inputs in their original form norms passes through
convolution and then through the complete form by \eqref{eq:window-support-1}.

Choose a nonnegative $\rho\in C_c^\infty(-1,1)$ with
$\int\rho=1$, and put $\rho_\eta(x)=\eta^{-1}\rho(x/\eta)$.
The multiplier satisfies $|\widehat{\rho_\eta}(\xi)|\le1$ and
converges pointwise to one.  Dominated convergence gives
\[
 \|\rho_\eta\star g-g\|_X\longrightarrow0.
\]
An explicit support-preserving form-core approximation is
\begin{equation}
 \theta_n=1-\frac1n,\qquad \eta_n=\frac{T}{2n},\qquad
 g^{[n]}=\rho_{\eta_n}\star U_{\theta_n}g\quad(n\ge2).
 \label{eq:window-support-5}
\end{equation}
Its support lies in
$[-T+T/(2n),T-T/(2n)]$, so $g^{[n]}\in C_c^\infty(I_T)$.
Moreover,
\[
 \|g^{[n]}-g\|_X
 \le\|U_{\theta_n}g-g\|_X
      +\|\rho_{\eta_n}\star g-g\|_X\longrightarrow0.
\]
Thus \eqref{eq:window-support-5} approximates endpoint-nonzero or discontinuous
admissible vectors without changing the original domain.

\subsubsection{Compact smooth tests and the exponential test algebra}

Let $\mathcal D(I)=C_c^\infty(I)$ carry its usual locally
convex inductive-limit topology.  On each fixed compact support
$K\Subset I$, its seminorms are
$\max_{j\le m}\|\partial^j f\|_\infty$.
This topology also gives a continuous inclusion into each
$\A_a$ of \eqref{eq:test-class}.  On a fixed compact set
$K\subset[-T,T]$,
\[
 p_{a,k}(f)\le 2T e^{aT}\|f^{(k)}\|_\infty.
\]
Thus compact smooth tests retain the convolution and involution of
the original exponential test algebra.  The embedding into the
logarithmic completion is continuous as well.
Indeed Fourier integration by parts twice gives, on each fixed
support $K$,
\[
 \|f\|_X\le C_K(\|f\|_\infty+\|f''\|_\infty),
\]
since the resulting Fourier square is bounded by a constant
times $(1+|\xi|)^{-4}$ and the weight grows logarithmically.
Zero extension along an open inclusion is continuous in this
topology.  For a fixed $g\in\mathcal D(I_T)$, the dilates
$U_\theta g$ converge to $g$ in $\mathcal D(I_T)$ as
$\theta\uparrow1$: their supports lie in one common compact
subset, and
\[
 \partial^jU_\theta g(x)
 =\theta^{-j-1/2}g^{(j)}(x/\theta)
\]
converges uniformly for every $j$.  Similarly, for small
$\eta$, $\rho_\eta\star g$ has one common compact support and
\begin{equation}
 \|\partial^j(\rho_\eta\star g-g)\|_\infty
 \le\sup_{|y|\le\eta}
         \|g^{(j)}(\cdot-y)-g^{(j)}\|_\infty\longrightarrow0.
 \label{eq:window-support-6}
\end{equation}
This is a direct approximation statement for the convolution;
it does not assert that $\rho_\eta$ converges to a function in
$\mathcal D$.

Convolution is jointly continuous in this inductive-limit
topology as well.  Here is a seminorm proof on $\mathbb R$.
Choose a locally finite smooth partition of unity $\chi_j$
with support in $[j-1,j+1]$.  Given a continuous seminorm $q$
on $\mathcal D(\mathbb R)$, choose $A_h\ge1$ and integers
$m_h$ so that
\[
 q(v)\le A_h\max_{r\le m_h}\|v^{(r)}\|_\infty
       \quad\text{if }\operatorname{supp}v\subset[h-2,h+2].
\]
Set
\[
 A'_j=\max_{|h|\le2|j|+2}A_h,\qquad
 m'_j=\max_{|h|\le2|j|+2}m_h,
\]
\[
 p(f)=\sum_{j\in\mathbb Z}A'_j
           \max_{r\le m'_j}\|\partial^r(\chi_jf)\|_\infty.
\]
This is a continuous seminorm: on each fixed compact-support
step only finitely many summands occur.  For
$f_j=\chi_jf$, $g_k=\chi_kg$, their convolution has support
in $[j+k-2,j+k+2]$.  If $|j|\ge|k|$, put the required
derivatives on $f_j$ and use $\|g_k\|_1\le2\|g_k\|_\infty$;
if $|k|>|j|$, reverse the roles.  Summing gives
\begin{equation}
 q(f\star g)\le2p(f)p(g).
 \label{eq:window-support-7}
\end{equation}
Thus convolution is continuous in the compact smooth test topology
and in the stated logarithmic completions.
For rough $g\in X_T$, the convergence in \eqref{eq:window-support-5} is a form-norm
statement; it is not a claimed convergence in $\mathcal D$ to a
nonsmooth limit.

\subsection{Effects of changing the complete form}
\label{sec:window-term-differences}

The following examples measure the effect of deleting a prime square,
replacing the polar term, or discarding the tail coupling.

Fix $T=3/4$ and $L=2T=3/2$.  The functions
\[
 f_0(x)=\mathbf1_{[-T,T]}(x),\qquad
 f_1(x)=x\mathbf1_{[-T,T]}(x)
\]
belong to $X_T$.  Their zero extensions are piecewise smooth with finitely many jumps.  Integration by parts
on the smooth pieces gives Fourier transforms $O(|\xi|^{-1})$; the
logarithmically weighted Fourier squares are integrable.
They are admissible without assigning zero endpoint values.
Formula \eqref{eq:window-support-5} supplies compact smooth versions in the same form
topology.

\subsubsection{The prime square with its von Mangoldt weight}

For $b=\log4$, the prescribed prime-square contribution is
\[
 -\frac{\log2}{2}\langle g,(S_b+S_{-b})g\rangle
 =-(\log2)\operatorname{Re} C_g(b),\qquad C_g(b)=\langle g,S_bg\rangle.
\]
Thus deleting $n=4$ changes the form by
$(\log2)\operatorname{Re} C_g(b)$; the weight is $\Lambda_{\mathrm{vM}}(4)=\log2$.
The two exact correlations are
\[
 C_{f_0}(b)=L-b,
\qquad
 C_{f_1}(b)=\frac{(L-b)((L-b)^2-3b^2)}{12}.
\]
Since $1<\log4<3/2$, the resulting changes have opposite,
strictly nonzero signs:
\begin{align*}
 \Delta_4(f_0)&=(\log2)(3/2-\log4)>0,\\
 \Delta_4(f_1)&=
 \frac{(\log2)(3/2-\log4)
       ((3/2-\log4)^2-3(\log4)^2)}{12}<0.
\end{align*}
In particular this deletion is not a positive perturbation on
all vectors.  The same formulas have the rational enclosures
\[
 \Delta_4(f_0)\in[0.07881474300351511479,
                         0.07881474300351511480],
\]
\[
 \Delta_4(f_1)\in[-0.03778186468138071530,
                         -0.03778186468138071529].
\]
These are differences between the specified terms, not negative
values of the complete form.

\subsubsection{The signed polar cross term}

Replacing
$2\operatorname{Re}(\overline{\ell_+}\ell_-)$ by
$|\ell_+|^2+|\ell_-|^2$ changes the form by exactly
\[
 \Delta_{\mathrm{pol}}(g)=|\ell_+(g)-\ell_-(g)|^2.
\]
The even vector $f_0$ annihilates this difference.  For the odd
vector $f_1$, $\ell_-(f_1)=-\ell_+(f_1)$ and
\[
 \ell_+(f_1)-\ell_-(f_1)
 =4\int_0^T x\sinh(x/2)\,dx
 =8\{T\cosh(T/2)-2\sinh(T/2)\}.
\]
Consequently
\begin{equation}
 \Delta_{\mathrm{pol}}(f_1)
 =64\{T\cosh(T/2)-2\sinh(T/2)\}^2
 >\frac{4T^6}{9}=\frac{81}{1024}>0.
 \label{eq:window-support-8}
\end{equation}
The strict lower bound uses $\sinh(x/2)>x/2$ for $x>0$.
The original polar value on $f_1$ is negative, whereas the
replacement is positive.  The polar and prime-square
differences persist for sufficiently large smooth
approximants \eqref{eq:window-support-5}, because their difference operators are
bounded on $L^2$.

\subsubsection{The nonzero omitted Schur correction}

Abbreviate $W_{3/4}$ to $W$.  Let $P=P_{64}^{3/4}$ be the Legendre projection onto degrees
$0,\ldots,63$, let $Q=1-P$, and let
$\mathcal K=QWQ$ denote the original complete tail operator.
Exact scalar inequalities give a positive lower bound for this tail.
The global kernel bound proved above gives
$\|\mathcal R_{3/4}\|\le3/8$.  In
\eqref{eq:window-arch-identity}, write
$c=-\gamma-\log(2\pi T)$ for the scalar term.

The elementary constants admit the following rational bounds.  The positive atanh series gives
\[
 \frac{69}{100}<\frac{56}{81}<\log2
       <\frac{25}{36}<\frac7{10},\qquad
 \log3<\frac{11}{10},\qquad \sqrt3>\frac{17}{10}.
\]
For the logarithm of three, four terms with a geometric
remainder suffice:
\[
 \log3\le2\sum_{j=0}^{3}\frac{(1/2)^{2j+1}}{2j+1}
       +\frac{2(1/2)^9}{9(1-1/4)}<\frac{11}{10}.
\]
The harmonic enclosure \eqref{eq:cofinal-gamma-trapezoid} applies.
At $q=2$, it follows that
$\gamma<41/32-69/100=473/800<3/5$.
The seven-term alternating upper bound for $\arctan1$ gives
\[
 \frac\pi4<1-\frac13+\frac15-\frac17
                    +\frac19-\frac1{11}+\frac1{13}<\frac56.
\]
The same alternating series with its next negative term gives
$\pi>3$.
Thus $3\pi/2<5$; and the first five positive exponential
terms at $17/10$ exceed five, so
$\log(3\pi/2)<\log5<17/10$.
Finally,
\[
 \nu(3/4)=\frac32\log2+\frac{\log3}{\sqrt3}
       <\frac{21}{20}+\frac{11}{17}<\frac74,
\]
is a norm bound for the complete prime operator.  Decomposition
into arithmetic chains for the shifts $\log2,\log3,\log4$
gives path-adjacency norms $\sqrt2,1,1$, respectively; multiplying
by their actual weights gives the displayed bound.  The signed
polar operator has eigenvalues $2T+2\sinh T$ and
$2T-2\sinh T$ on the spans of its even and odd eigenvectors.  Thus its
negative part is bounded using
\[
 2(\sinh(3/4)-3/4)
 \le\frac{(3/4)^3}{3(1-(3/4)^2/20)}
 =\frac{45}{311}<\frac15.
\]
Since $V_T\ge0$, these inequalities prove
\[
 W\succeq
 D_{3/4}-\left(\frac35+\frac{17}{10}+\frac38+\frac74+\frac15\right)I
 =D_{3/4}-\frac{37}{8}I.
\]
Exact rational summation gives
\[
 H_{64}=
 \frac{623171679694215690971693339}
      {131362987122535807501262400}>\frac{47}{10};
\]
therefore
\[
 \mathcal K\succeq(H_{64}-37/8)Q>\frac3{40}Q.
\]
This complement is positive regardless of the sign of the finite low block.
For $u\in P\mathcal H_{3/4}$,
the omitted correction is
\[
 \mathcal B(u)=\langle QWu,\mathcal K^{-1}QWu\rangle,
\]
so the exact Schur form is $Q_{3/4}(u)-\mathcal B(u)$.

This correction is nonzero already on $u=1$.  Indeed
\[
 W1=V_T+h_T,\qquad h_T\in L^\infty(I_T),\qquad
 V_T=-\tfrac12\log(1-(x/T)^2).
\]
The regular convolution, polar term and every translated
constant are bounded.  The essential endpoint logarithmic
divergence prevents $W1$ from being any finite polynomial.
Equivalently, on a sufficiently small one-sided neighborhood
of $T$ that contains no translation cutoff,
\[
 (W1)'(x)=\frac{x}{T^2-x^2}+O(1)
       =\frac1{2(T-x)}+O(1).
\]
Therefore $QW1\ne0$, and positivity and injectivity of
$\mathcal K^{-1}$ imply $\mathcal B(1)>0$.

A single Legendre coefficient gives a quantitative lower bound.  Write
$p_n(x)=P_n(x/T)$ and $g_n=\|p_n\|_2^2=L/(2n+1)$, and
take $h=p_{66}\in Q\mathcal Q_{3/4}$.  Define
\[
 \alpha_m=1-\frac{\log m}{T}\quad(m=2,3,4),
\]
\begin{equation}
 c_{66}^{\mathrm{alg}}=
 \frac{L}{66\cdot67}
 -\frac{L}{133}\sum_{m=2,3,4}
       \frac{\Lambda_{\mathrm{vM}}(m)}{\sqrt m}
       \{P_{67}(\alpha_m)-P_{65}(\alpha_m)\}.
 \label{eq:window-support-9}
\end{equation}
This retains all prime powers with their prescribed weights.
The first term is the exact logarithmic-potential matrix
entry.  The translated-constant integrals follow from
$\int P_n=(P_{n+1}-P_{n-1})/(2n+1)$ and parity.

The degree-$64$ polynomial regular kernel maps $1$ to a
polynomial of degree at most $65$, and the degree-$32$ polar
operator maps it to degree at most $32$.  Orthogonality
annihilates both against $p_{66}$.  The requisite approximation
bound is analytic: on $|z|=3$,
\[
 |\sinh z|\ge\sin3,\qquad
 |R(z)|\le\frac{e^{3/2}}{2\sin3}+\frac16<17.
\]
For the first inequality, write $z=x+iy$ with $x^2+y^2=9$.
The derivative of $\sinh^2x+\sin^2y$ with respect to $x^2$
is $\sinh(2|x|)/(2|x|)-\sin(2|y|)/(2|y|)\ge0$, with the
continuous values at zero.  Its minimum is therefore $\sin^23$.
The last inequality follows already from the elementary Taylor
bounds $e^{3/2}<9/2$ and $\sin3>7/50$.  Cauchy's estimate and
the integral-kernel norm bound then give
\[
 \|\mathcal R_T-\mathcal R_{T,64}\|
 \le\frac{51}{2}\,2^{-64}.
\]
For the degree-$32$ polar Taylor approximation, putting
$\eta=(3/2)(3/8)^{33}/33!$ gives operator error at most
$(13/2)\eta+(25/8)\eta^2$.  The sum of these two explicit
rational bounds is less than $2\cdot10^{-18}$.  Together with
$\|p_{66}\|_2\|1\|_2=L/\sqrt{133}<1$, this gives
\begin{equation}
 |\langle W1,p_{66}\rangle_2-c_{66}^{\mathrm{alg}}|<2\cdot10^{-18}.
 \label{eq:window-support-10}
\end{equation}
No replacement of Euler's constant is needed in this off-diagonal
entry, since its scalar contribution is exactly orthogonal.

To enclose this coefficient by rational arithmetic, use
\[
 \log z=2\sum_{j=0}^{s-1}\frac{q^{2j+1}}{2j+1}+\epsilon_s,
 \qquad q=\frac{z-1}{z+1},\quad
 0\le\epsilon_s\le\frac{2q^{2s+1}}{(2s+1)(1-q^2)}
 \quad(z>1).
\]
At $z=2,3$ take $s=240$, and set $\log4=2\log2$.
Enclose $\sqrt2,\sqrt3$ by rational numbers whose squares
bracket $2,3$.  The recurrence
\[
 P_0(t)=1,\quad P_1(t)=t,\quad
 P_{n+1}(t)=\frac{(2n+1)tP_n(t)-nP_{n-1}(t)}{n+1}
\]
then encloses every term in \eqref{eq:window-support-9}.  Combining it with \eqref{eq:window-support-10} gives
\begin{equation}
 \langle W1,p_{66}\rangle_2
 \in[-0.0010128911888997950,
       -0.0010128911888997909].
 \label{eq:window-support-11}
\end{equation}
This calculation involves one coefficient of the original operator.
Its error includes the regular and polar remainders.

For the denominator, the diagonal logarithmic identity is
\[
 (V_T)_{nn}=2(H_{2n}-H_n)+\frac1{2n+1}-\log2
 <\log2+\frac1{2n+1}\quad(n\ge1).
\]
It follows by differentiating the monic Jacobi squared norm
with weight $(1-s^2)^\alpha$ at $\alpha=0$; the derivative
of the monic polynomial is orthogonal to that polynomial.
The inequality uses $H_{2n}-H_n<\log2$. Thus $H_{66}<5$,
the normalized logarithmic diagonal entry is less than one, and
$-\gamma-\log(2\pi T)<0$.  The kernel norm bound proved above
is $\|\mathcal R_T\|\le3/8$, the polar norm is less
than four, and the complete prime-operator norm satisfies
\[
 \nu(3/4)=\frac32\log2+\frac{\log3}{\sqrt3}<2.
\]
The last formula includes the $n=4$ term.  Hence
\begin{equation}
 0<Q_{3/4}(p_{66})<13g_{66}=\frac{39}{266}.
 \label{eq:window-support-12}
\end{equation}
The positive lower inequality follows from the original tail
comparison.  The upper inequality includes every term of the
same complete form.

The variational identity for the exact inverse gives
\[
 \mathcal B(1)
 =\sup_{v\in Q\mathcal Q_{3/4}}
    \{2\operatorname{Re}\langle W1,v\rangle_2-Q_{3/4}(v)\}
 \ge\frac{|\langle W1,p_{66}\rangle_2|^2}{Q_{3/4}(p_{66})}.
\]
In particular \eqref{eq:window-support-11}--\eqref{eq:window-support-12} imply the simple exact lower bound
\[
 \mathcal B(1)>\frac{266}{39\cdot10^6}>0.
\]
Keeping the interval endpoints gives the sharper bounds
\begin{equation}
 \mathcal B(1)>6.99749531042362\cdot10^{-6},\qquad
 \mathcal B\bigl(1/\sqrt L\bigr)
       >4.66499687361574\cdot10^{-6}.
 \label{eq:window-support-13}
\end{equation}
Thus omitting the complete Schur correction changes an actual
low-space value by a strictly positive amount.  This bound concerns the original inverse $\mathcal K^{-1}$.

If the alteration is implemented by deleting the off-diagonal
$P$--$Q$ block of the full operator, its effect can also be
seen directly on $1+p_{66}$: the change in its full quadratic
value is
\[
 \Delta_{\mathrm{block}}(1+p_{66})
 =-2\langle W1,p_{66}\rangle_2>\frac1{500}.
\]
Since $\|1+p_{66}\|_2^2=L+g_{66}=201/133$, the change per
unit physical norm is greater than $133/100500$.
The cross block is bounded, because $QWP$ has finite-dimensional
domain and polynomial columns in the operator domain.  Hence
this changed-form difference also persists for compact smooth
approximants \eqref{eq:window-support-5}.  This distinguishes block deletion from
merely omitting a term in the finite Schur formula; neither
calculation asserts that the original value $Q_{3/4}(1+p_{66})$ is negative.

